

\begin{table}
\centering
\begin{tabular}{| c | c |}
\hline
\# of observed events	                                          	& P-value    		    	\\
\hline
$N_{obs} = 2$      		                    	& 0.525                    	\\
$ N_{obs} = 3$      		                    	& 0.264                     	\\
$ \bf N_{obs} = 4$      		                 & \bf 0.109                      	\\
$N_{obs} = 5$      		                    	& 0.0352                      	\\
$N_{obs} = 6$      		                    	& 0.0108                     	\\
\hline
\end{tabular}
\caption{
Summary of the expected rate-only p-values for different number of 
observed events. The simulated number of POT is 
\TotalPot POT.
%\FakePotLow.
\NumToysLarge toys are used to get the distribution of the test statistic and the 
Asimov data is used to compute the p-value.
The Asimov dataset for $\beta=1$ is in bold.
}
\label{tab:summary_sensi_rateonly_highpot}
\end{table}



%\begin{table}
%\centering
%\begin{tabular}{| c | c |}
%\hline
%\# of observed events	                                          	& P-value    		    	\\
%\hline
%$N_{obs} = 2$      		                    	& 0.595                      	\\
%$N_{obs} = 3$      		                    	& 0.322                      	\\
%$N_{obs} = 4$      		                    	& 0.144                    	\\
%$\bf N_{obs} = 5$      		                    	& \bf 0.0538                      	\\
%$N_{obs} = 6$      		                    	& 0.0174                     	\\
%\hline
%\end{tabular}
%\caption{
%Summary of the expected rate-only p-values for different number of 
%observed events. The simulated number of POT is \FakePotHigh.
%\NumToys toys are used to get the distribution of the test statistic and the 
%Asimov data is used to compute the p-value.
%The Asimov dataset for $\beta=1$ is in bold.
%}
%\label{tab:summary_sensi_rateonly_lowpot}
%\end{table}


\begin{table}
\centering
\begin{tabular}{| c | c | c |}
\hline
P-value	                                          	& Asimov    	& Mean 		    	\\
\hline
Rate-only                          			& 0.109           	& 0.203                  	\\
Rate+shape                          		& 0.0699          	& 0.139                  	\\
\hline
\end{tabular}
\caption{
Summary of the expected rate-only and rate+shape p-values for 
\TotalPot POT.
%\FakePotLow POT.
The Asimov p-value is calculated using the Asimov dataset while the 
mean p-value is calculated making \NumToysLarge toys and taking the mean 
of the p-value distribution.
}
\label{tab:summary_sensi_lowpot}
\end{table}

%\begin{table}
%\centering
%\begin{tabular}{| c | c | c |}
%\hline
%P-value	                                          	& Asimov    	& Mean 		    	\\
%\hline
%Rate-only                          			& 0.0538           & 0.188                   	\\
%Rate+shape                          		& 0.0572          	& 0.127                  	\\
%\hline
%\end{tabular}
%\caption{
%Summary of the expected rate-only and rate+shape p-values for \FakePotHigh POT.
%The Asimov p-value is calculated using the Asimov dataset while the 
%mean p-value is calculated making \NumToys toys and taking the mean 
%of the p-value distribution.
%}
%\label{tab:summary_sensi_highpot}
%\end{table}
%






% Comparison with many toys

\begin{table}
\centering
\begin{tabular}{| c | c | c |}
\hline
P-value	                                          	& Asimov    	& Mean 		    	\\
\hline
Rate-only (\NumToys toys)         			& 0.110           	& 0.207                  	\\
Rate+shape (\NumToys toys)              		& 0.0726          	& 0.143                  	\\
\hline
Rate-only (\NumToysLarge toys)        			& 0.109           & 0.203                	\\
Rate+shape (\NumToysLarge toys)             		& 0.0699          	& 0.139                	\\
\hline
\end{tabular}
\caption{
Summary of the expected rate-only and rate+shape p-values for 
\TotalPot POT, using \NumToys (up) and \NumToysLarge (bottom) toys for fake datasets (both $\beta=0$ and $\beta=1$).
The Asimov p-value is calculated using the Asimov dataset while the mean one
from \NumToys (up) and \NumToysLarge (bottom) toys and taking the mean 
of the p-value distribution.
}
\label{tab:summary_sensi_manytoys}
\end{table}


\begin{table}
\centering
\begin{tabular}{| c | c | c |}
\hline
P-value	                                          	& Asimov    	& Mean 		    	\\
\hline
Rate-only (\NumToys toys)        			& 0.0343           & 0.1004                   	\\
Rate+shape (\NumToys toys)             		& 0.0249          	& 0.0649                 	\\
\hline
Rate-only (\NumToysLarge toys)        			& 0.0331           & 0.1007                   	\\
Rate+shape (\NumToysLarge toys)             		& 0.0243          	& 0.0660                 	\\
\hline
\end{tabular}
\caption{
Summary of the expected rate-only and rate+shape p-values for 
$8 \times 10^{20}$ POT, using \NumToys (up) and \NumToysLarge (bottom) toys for fake datasets (both $\beta=0$ and $\beta=1$).
The Asimov p-value is calculated using the Asimov dataset while the mean one
from \NumToys (up) and \NumToysLarge (bottom) toys and taking the mean 
of the p-value distribution.
}
\label{tab:summary_sensi_manytoys}
\end{table}

%





\begin{table}
\centering
\begin{tabular}{| c | c | c |}
\hline
		                                          	& p-value $\leq 0.10$    		& p-value $\leq 0.05$ 		\\
\hline
Fraction of toys	(rate-only)			& 44.3\%					& 44.3\%				\\
Fraction of toys	(rate+shape) 		& 64.0\%					& 52.8\%				\\
\hline
\end{tabular}
\caption{
The probability to have a p-value smaller than 0.05 and 0.10 is shown.
\NumToysLarge toys for
%\FakePotLow POT 
\TotalPot POT
are produced.
Then for each toy the p-value is computed as described in section~\ref{subsec:pvalue_mean}.
}
\label{tab:pvalue_mean_0p05_lowpot}
\end{table}

%
%\begin{table}
%\centering
%\begin{tabular}{| c | c | c |}
%\hline
%		                                          	& p-value $\leq 0.10$    		& p-value $\leq 0.05$     	\\
%\hline
%%Mean p-value (rate-only) 			&  0.0213     				&  0.00809                  	\\
%Fraction of toys	(rate-only)			& 53.0\%					&  37.7\%				\\
%%Mean p-value (rate+shape)       		&  0.0214    				&  0.0118                  	\\
%Fraction of toys	(rate+shape) 		& 67.1\%					&  56.6\%				\\
%\hline
%\end{tabular}
%\caption{
%The probability to have a p-value smaller than 0.05 and 0.10 is shown.
%\NumToys toys for \FakePotHigh POT are produced.
%Then for each toy the p-value is computed as described in~\ref{subsec:pvalue_mean}.
%}
%\label{tab:pvalue_mean_0p05_highpot}
%\end{table}
%
%


\begin{table}
\centering
\begin{tabular}{| c | c | c |}
\hline
\# obs events $>$ \# exp bkg		& p-value    	 & Fraction of toys    	\\
\hline
Rate-only                          			&   0.149    	& 92.4\%             	\\
Rate+shape                          		&   0.0927   	& 92.4\%                 	\\
\hline
\end{tabular}
\caption{
Summary of the expected rate-only and rate+shape mean p-values for 
\TotalPot POT.
%\FakePotLow POT.
The Asimov p-value is calculated using the Asimov dataset while the 
mean p-value is calculated making \NumToysLarge toys and taking the mean 
of the p-value distribution.
Only fake datasets with number of observed events higher than the
expected background are considered.
The fraction of selected toys is shown as well. 
}
\label{tab:pvalue_mean_bkg_lowpot}
\end{table}

%
%\begin{table}
%\centering
%\begin{tabular}{| c | c | c |}
%\hline
% \# obs events $>$ \# exp bkg           & P-value  	 & Fraction of toys    	\\
%\hline
%Rate-only                          			&  0.151      	&   95\%               	\\
%Rate+shape                          		&  0.0947    	&   95\%             	\\
%\hline
%\end{tabular}
%\caption{
%Summary of the expected rate-only and rate+shape mean p-values for \FakePotHigh POT.
%The Asimov p-value is calculated using the Asimov dataset while the 
%mean p-value is calculated making \NumToys toys and taking the mean 
%of the p-value distribution.
%Only fake datasets with number of observed events higher than the
%expected background are considered.
%The fraction of selected toys is shown as well. 
%}
%\label{tab:pvalue_mean_bkg_highpot}
%\end{table}


\begin{table}
\centering
\begin{tabular}{| c | c |}
\hline
p-value (rate-only) $>$ p-value (rate+shape)		& Probability          	\\
\hline
All toys									& 85.0\%									\\
Toys w/ \# obs. events $>$ \# exp. bkg 			& 77.4\%									\\
\hline
\end{tabular}
\caption{
The probability to have a rate-only p-value larger than the rate+shape p-value is shown.
\NumToysLarge toys for 
\TotalPot POT.
%\FakePotLow POT are produced.
Then for each toy the p-value is computed as described in section~\ref{subsec:pvalue_mean}.
Both all toys are considered as well as the toys with the number of observed events larger
than the number of expected background.
}
\label{tab:pvalue_mean_rate_large_shape_lowpot}
\end{table}


%\begin{table}
%\centering
%\begin{tabular}{| c | c |}
%\hline
%p-value (rate-only) $>$ p-value (rate+shape)		& Probability          	\\
%\hline
%All toys									& 81.4\%									\\
%Toys w/ \# obs. events $>$ \# exp. bkg 			& 77.4\%									\\
%\hline
%\end{tabular}
%\caption{
%The probability to have a rate-only p-value larger than the rate+shape p-value is shown.
%\NumToys toys for \FakePotHigh POT are produced.
%Then for each toy the p-value is computed as described in~\ref{subsec:pvalue_mean}.
%Both all toys are considered as well as the toys with the number of observed events larger
%than the number of expected background.
%}
%\label{tab:pvalue_mean_rate_large_shape_highpot}
%\end{table}

